% runtime/runtime_services.tex

\chapter{Runtime Services}

\label{chap:runtime-services}

\index{Runtime services}
\index{Architecture!runtime}

Runtime services are the execution layer that evaluates approved
engineering knowledge in operational contexts. They consume established
objects from Foundations and State and generate derived runtime
quantities. They are not engineering objects themselves and do not
redefine upstream semantics.

This chapter answers how runtime evaluation is organized so that
projection, frame construction, dynamic quantities, and representation
boundaries remain consistent with canonical AIM semantics. It
establishes the runtime service architecture and its engineering
consequences for downstream interpretation.

% ============================================================
\section{Runtime Interpretation}
% ============================================================

\begin{definition}[Runtime service]
\label{def:runtime-service}
\index{Runtime service!definition}
A runtime service is an operator-based evaluation process acting on
alignment states during operational use.
\end{definition}

\begin{principle}[Runtime evaluation principle]
\label{princ:runtime-evaluation}
Operational quantities should be evaluated from underlying state
descriptions rather than stored redundantly.
\end{principle}

Runtime services evaluate geometry, projections, frames,
realization-dependent states, dynamic quantities, and vehicle-space
interpretation through explicit operator composition.

The remaining sections now follow the same reader path: from runtime
state evaluation, to concrete services, to pipeline composition, and
finally to representation boundaries and architectural consequence.

% ============================================================
\section{Runtime State Evaluation}
% ============================================================

Runtime evaluation is governed by the canonical evaluation operator in
\cref{def:evaluation-operator}.

\begin{definition}[Runtime evaluation]
\label{def:runtime-evaluation}
\index{Runtime evaluation}
Runtime evaluation is represented by
\[
\mathcal E
:
\mathcal X
\times
\Omega_{\mathrm{track}}
\rightarrow
\mathcal Y.
\]
\end{definition}

Here \(\mathcal X\) denotes railway state space,
\(\Omega_{\mathrm{track}}\) intrinsic track space, and
\(\mathcal Y\) runtime-evaluated quantities.

Typical runtime outputs include world coordinates, tangent and frame
quantities, curvature/cant-derived quantities, projection states, and
vehicle-space-derived quantities.

% ============================================================
\section{Persistent Information versus Derived Runtime State}
% ============================================================

Within AIM, persistent data stores sparse, semantically stable
engineering information. Runtime services reconstruct derived quantities
from that information.

\begin{principle}[Runtime reconstruction principle]
\label{princ:runtime-reconstruction}
\index{Runtime reconstruction principle}
Runtime geometry should be reconstructed from sparse alignment semantics
rather than stored redundantly.
\end{principle}

This separation avoids synchronization instability, redundant storage,
and realization inconsistencies.

\begin{figure}[ht]
\centering
% figures/runtime/support_runtime_to_representation_boundary.tex
\begin{tikzpicture}[
  node distance=1.4cm,
  box/.style={draw, rounded corners, minimum width=4.7cm, minimum height=0.85cm, align=center},
  sec/.style={draw, rounded corners, dashed, minimum width=4.7cm, minimum height=0.85cm, align=center},
  arrow/.style={->, thick}
]
  \node[box] (rt) {Runtime};
  \node[sec, below=of rt] (rep) {Representation};
  \draw[arrow] (rt) -- (rep);
  \draw[thin] ($(rt.center)+(-3.0,-0.95)$) -- ($(rt.center)+(3.0,-0.95)$);
\end{tikzpicture}

\caption{Supporting boundary: runtime evaluation remains upstream of representation.}
\label{fig:support-runtime-representation-boundary}
\end{figure}

The figure makes the ownership boundary explicit: runtime produces
derived engineering state, while representation consumes that state for
encoding or visualization. The engineering consequence is that runtime
services must be validated against evaluation correctness, not against
format-specific rendering choices.

% ============================================================
\section{Runtime pose3 Evaluation}
% ============================================================

The foundational pose3 definition is established in
\cref{def:pose3,def:runtime-pose3}. Runtime services consume these
concepts and evaluate them operationally.

\begin{definition}[Runtime pose3 evaluation]
\label{def:runtime-pose3-evaluation}
\index{State!pose3!runtime evaluation}
Runtime pose3 evaluation may be represented schematically as
\[
\mathrm{pose3}_{\mathrm{run}}(s)
=
\mathcal E_{\mathrm{pose3}}(x,s).
\]
\end{definition}

The runtime pose3 state is therefore a derived evaluation state, not a
persistent engineering identity.

% ============================================================
\section{Projection Services}
% ============================================================

Runtime projection services consume the canonical operators in
\cref{def:track-to-world-operator,def:world-to-track-operator}.

The next two figures answer a directional question before the service
definitions are stated: which operator maps from world to track, and
which provides the complementary mapping from track to world?

\begin{figure}[ht]
\centering
% figures/runtime/support_world_to_track_operator.tex
\begin{tikzpicture}[
  node distance=1.25cm,
  box/.style={draw, rounded corners, minimum width=3.4cm, minimum height=0.8cm, align=center},
  op/.style={draw, rounded corners, minimum width=3.9cm, minimum height=0.85cm, align=center},
  arrow/.style={->, thick}
]
  \node[box] (world) {world};
  \node[op, below=of world] (opw) {world-to-track operator};
  \node[box, below=of opw] (track) {track};
  \draw[arrow] (world) -- (opw);
  \draw[arrow] (opw) -- (track);
\end{tikzpicture}

\caption{Supporting operator view: world-to-track mapping as a directed runtime operator.}
\label{fig:support-world-to-track-operator}
\end{figure}

\begin{figure}[ht]
\centering
% figures/runtime/support_track_to_world_operator.tex
\begin{tikzpicture}[
  node distance=1.25cm,
  box/.style={draw, rounded corners, minimum width=3.4cm, minimum height=0.8cm, align=center},
  op/.style={draw, rounded corners, minimum width=3.9cm, minimum height=0.85cm, align=center},
  arrow/.style={->, thick}
]
  \node[box] (track) {track};
  \node[op, below=of track] (opt) {track-to-world operator};
  \node[box, below=of opt] (world) {world};
  \draw[arrow] (track) -- (opt);
  \draw[arrow] (opt) -- (world);
\end{tikzpicture}

\caption{Supporting operator view: track-to-world mapping as the complementary directed runtime operator.}
\label{fig:support-track-to-world-operator}
\end{figure}

Read together, the pair establishes a directional runtime contract
before formal service definitions: inverse projection and forward
projection are related but non-identical operators. The engineering
consequence is explicit separation of numerical strategy, ambiguity
handling, and quality criteria for each mapping direction.

\begin{definition}[Track-to-world runtime service]
\label{def:track-to-world-runtime}
\index{Projection!track-to-world runtime service}
The track-to-world runtime service evaluates
\[
\mathcal P_{\mathrm{tw}}
:
\Omega_{\mathrm{track}}
\rightarrow
\Omega_{\mathrm{world}}.
\]
\end{definition}

\begin{definition}[World-to-track runtime service]
\label{def:world-to-track-runtime}
\index{Projection!world-to-track runtime service}
The inverse runtime projection service evaluates
\[
\mathcal P_{\mathrm{wt}}
:
\Omega_{\mathrm{world}}
\rightarrow
\Omega_{\mathrm{track}}.
\]
\end{definition}

The inverse service is nonlinear, potentially ambiguous,
metric-dependent, and context-dependent. Projection remains distinct
from CRS transformation, as established in
\cref{def:coordinate-transform-operator}.

This distinction is the practical boundary condition for all later
runtime services: projection interprets intrinsic geometry, while CRS
transformation only changes world-space representation.

% ============================================================
\section{Frame Services}
% ============================================================

Runtime frame construction consumes the canonical frame operator in
\cref{def:frame-operator}.

\begin{definition}[Frame service]
\label{def:frame-service}
\index{Reference frame!runtime service}
The frame service evaluates
\[
\mathcal F(x_{\mathrm{metric}},s)
=
(\mathbf t(s),\mathbf n(s),\mathbf b(s)).
\]
\end{definition}

Frame services support cant interpretation, pose3 evaluation,
vehicle-space interpretation, projection interpretation, and runtime
dynamics. Their evaluation may depend on metric realization.

% ============================================================
\section{Derived Dynamic Runtime Quantities}
% ============================================================

Many operational quantities are derived at runtime and are not
persistent alignment knowledge.

\begin{definition}[Effective lateral acceleration]
\label{eq:runtime-effective-acceleration}
\index{Acceleration!effective lateral}
\[
a_{\mathrm{eff}}
=
v^2\kappa-g\sin\phi.
\]
\end{definition}

This is the runtime form of the canonical relation in
\cref{eq:effective-lateral-acceleration}.

\begin{definition}[Runtime jerk]
\label{eq:runtime-jerk}
\index{Jerk!runtime}
\[
j
=
v^3\kappa'
-
gv\cos\phi\,\phi'.
\]
\end{definition}

This is the runtime form of the canonical jerk relation in
\cref{eq:jerk-relation}.

% ============================================================
\section{Hybrid and LOD-Dependent Evaluation}
% ============================================================

\begin{principle}[Hybrid runtime principle]
\label{princ:hybrid-runtime}
Efficient runtime systems combine approximate coarse evaluation, cached
representations, and exact local refinement.
\end{principle}

\begin{definition}[LOD-dependent evaluation]
\label{def:lod-evaluation}
\index{Level of detail (LOD)!runtime evaluation}
Runtime evaluation depends on scale, precision requirements, visibility,
operational context, interaction state, and runtime cost.
\end{definition}

This combination is central for scalable projection, frame evaluation,
visualization, and vehicle-space interpretation.

% ============================================================
\section{Caching and Consistency}
% ============================================================

\begin{definition}[Runtime cache]
\label{def:runtime-cache}
\index{Runtime cache}
A runtime cache stores previously evaluated quantities to avoid repeated
computation.
\end{definition}

Caching can accelerate projection, frame, sampled geometry, and dynamic
queries, while introducing consistency and invalidation requirements.

% ============================================================
\section{Metric-Aware Runtime}
% ============================================================

Runtime evaluation depends on metric realization
(\cref{chap:metric-realization}). Different realizations affect distance
interpretation, curvature-related quantities, frame construction,
projection behaviour, and stationing interpretation.

\begin{principle}[Metric-aware runtime principle]
\label{princ:metric-aware-runtime}
Runtime systems should evaluate geometry through metric-aware operator
services rather than through hardcoded Euclidean assumptions.
\end{principle}

% ============================================================
\section{Runtime Pipelines and Service Interaction}
% ============================================================

\begin{definition}[Runtime operator pipeline]
\label{def:runtime-pipeline}
\index{Pipeline!runtime operator}
A runtime pipeline is a composed operator chain such as
\[
\mathcal E_{\mathrm{run}}
\sim
\mathcal P
\circ
\mathcal F
\circ
\mathcal G.
\]
\end{definition}

The following figures are read as progressively richer answers to the
same question: how does the operator chain change when realized geometry
must be included explicitly?

\begin{figure}[ht]
\centering
% figures/runtime/runtime_pipeline.tex

\begin{tikzpicture}[
  node distance=0.85cm,
  box/.style={draw, rounded corners, align=center, minimum width=5.0cm, minimum height=0.7cm},
  arrow/.style={->, thick}
]
\node[box] (sparse) {Sparse Alignment};
\node[box, below=of sparse] (pose2) {pose2 Reconstruction};
\node[box, below=of pose2] (pose3) {pose3 Runtime State};
\node[box, below=of pose3] (dyn) {Runtime Dynamics\\$\mathcal D$};
\node[box, below=of dyn] (vehicle) {Vehicle Interpretation\\$\mathcal V$};
\node[box, below=of vehicle] (opt) {Operational Evaluation\\$\mathcal O$};
\draw[arrow] (sparse) -- (pose2);
\draw[arrow] (pose2) -- (pose3);
\draw[arrow] (pose3) -- (dyn);
\draw[arrow] (dyn) -- (vehicle);
\draw[arrow] (vehicle) -- (opt);
\end{tikzpicture}

\caption{Conceptual runtime evaluation pipeline.}
\label{fig:runtime-pipeline}
\end{figure}

This figure is interpreted as the baseline operator composition used for
target-state runtime evaluation.

\begin{definition}[Realized runtime pipeline]
\label{def:realized-runtime-pipeline}
\index{Pipeline!realization-aware runtime}
When realized geometry is required, runtime evaluation may conceptually
depend on
\[
\mathcal E_{\mathrm{real}}
\sim
\mathcal P
\circ
\mathcal F
\circ
\mathcal G
\circ
\mathcal R.
\]
\end{definition}

\begin{figure}[ht]
\centering
% figures/runtime/realization_pipeline.tex

\begin{tikzpicture}[
  node distance=0.9cm,
  box/.style={draw, rounded corners, align=center, minimum width=5.0cm, minimum height=0.75cm},
  arrow/.style={->, thick}
]
\node[box] (target) {Target State\\$\mathrm{pose3}_{\mathrm{target}}$};
\node[box, below=of target] (realization) {Realization Operator\\$\mathcal R_{\mathrm{pose3}}$};
\node[box, below=of realization] (real) {Realized State\\$\mathrm{pose3}_{\mathrm{real}}$};
\node[box, below=of real] (measured) {Measured / Reconstructed State};
\node[box, below=of measured] (runtime) {Runtime Operation};
\draw[arrow] (target) -- (realization);
\draw[arrow] (realization) -- (real);
\draw[arrow] (real) -- (measured);
\draw[arrow] (measured) -- (runtime);
\end{tikzpicture}

\caption{Realization-aware runtime evaluation pipeline.}
\label{fig:realization-pipeline}
\end{figure}

In contrast, the realization-aware figure adds \(\mathcal R\) explicitly
to represent realized-state dependence. The engineering consequence is
that service selection must declare whether target or realized state is
required, because this choice changes the active operator chain.

Different runtime services activate different operator chains, depending
on precision, context, and purpose.

This provides the transition to the final boundary statement: runtime
evaluation may vary operationally, but it must not alter conceptual
ownership of engineering meaning.

% ============================================================
\section{Runtime, Realization, and Representation Boundaries}
% ============================================================

Runtime must distinguish target, realized, measured, and temporary
interaction states. Runtime evaluation and physical realization are
coupled but conceptually distinct: realization transforms target states,
while runtime evaluates states.

Representation remains downstream. Runtime services provide derived
state and quantities; representation services display, encode, or
exchange those results for specific applications.

% ============================================================
\section{Conceptual Runtime Architecture}
% ============================================================

\begin{figure}[ht]
\centering

\begin{tikzpicture}[
node distance=1.5cm and 2cm,
box/.style={
draw,
rounded corners,
align=center,
minimum width=3.6cm,
minimum height=1.0cm
},
arrow/.style={->, thick}
]

\node[box] (sparse) {
Sparse Alignment\\
Semantic Representation
};

\node[box, below=of sparse] (metric) {
Metric Realization\\
\(\mathcal G\)
};

\node[box, below=of metric] (eval) {
Runtime Evaluation\\
\(\mathcal E\)
};

\node[box, below left=of eval] (proj) {
Projection Services\\
\(\mathcal P\)
};

\node[box, below=of eval] (frame) {
Frame Services\\
\(\mathcal F\)
};

\node[box, below right=of eval] (dyn) {
Dynamic Evaluation
};

\node[box, below=2.2cm of eval] (vehicle) {
Derived Runtime State\\
Application Interpretation
};

\draw[arrow] (sparse) -- (metric);
\draw[arrow] (metric) -- (eval);

\draw[arrow] (eval) -- (proj);
\draw[arrow] (eval) -- (frame);
\draw[arrow] (eval) -- (dyn);

\draw[arrow] (proj) -- (vehicle);
\draw[arrow] (frame) -- (vehicle);
\draw[arrow] (dyn) -- (vehicle);

\end{tikzpicture}

\caption{
Runtime services transform persistent engineering information through
metric-aware evaluation into derived runtime state for downstream
interpretation.
}
\label{fig:runtime-architecture}
\index{Figures!runtime architecture}

\end{figure}

This architecture figure summarizes the chapter argument in executable
form: persistent semantics feed metric-aware evaluation, and evaluation
fans out into service families that converge into derived runtime state.
The engineering consequence is a stable integration contract for
downstream Reality, Modelling, and Optimization chapters.

% ============================================================
\section{Canonical Interpretation}
% ============================================================

\begin{proposition}
Within AIM, runtime services are operator-based evaluation systems
acting on established engineering and state concepts.
\end{proposition}

\begin{proposition}
Derived runtime quantities are evaluation results and are not persistent
engineering identity.
\end{proposition}

\begin{proposition}
Projection, frame construction, metric realization, and runtime dynamics
form a coupled service architecture.
\end{proposition}

\begin{proposition}
Representation remains downstream of runtime evaluation.
\end{proposition}

% ============================================================
\section{Connection to AIM}
% ============================================================

\begin{summary}
AIM runtime services consume persistent engineering information and
state-layer concepts and generate derived runtime engineering state
through operator-based evaluation.

This architecture keeps semantics stable, keeps evaluation explicit,
maintains representation boundaries, and enables
application-specific interpretation without redefining upstream
knowledge.

Therefore, the runtime part closes with a precise consequence for the
next parts: downstream reality, modelling, optimization, and
applications may specialize runtime services, but they do so by
consuming this architecture rather than redefining it.
\end{summary}
